\documentclass[11pt]{article}

\usepackage[T1]{fontenc}
\usepackage{amsmath,amssymb,amsthm}
\usepackage{enumitem}
\usepackage[margin=1in]{geometry}
\usepackage[colorlinks=true,linkcolor=blue,citecolor=blue,urlcolor=blue]{hyperref}

\newtheorem{theorem}{Theorem}
\newtheorem{lemma}{Lemma}
\newtheorem{proposition}{Proposition}
\newtheorem{remark}{Remark}

\newcommand{\Lip}{\operatorname{Lip}}
\newcommand{\NS}{\mathsf{NS_T}}

\title{A Non-Ultrametric Example with Positive Snowflake Self-Embedding Parameter}
\author{Agent lane 06}
\date{2026-07-18}

\begin{document}
\maketitle

\section*{Source Metadata}

This packet concerns arXiv:1206.3774, \emph{Embeddability of snowflaked metrics
with applications to the nonlinear geometry of the spaces \(L_p\) and
\(\ell_p\) for \(0<p<\infty\)}, by Fernando Albiac and Florent Baudier.  The
arXiv page records the paper as submitted on 2012-06-17, revised on
2012-10-02, and later published in \emph{Journal of Geometric Analysis} 25
(2015), no. 1, 1--24.

The relevant open question appears near the end of the paper as Question
\(\mathrm{q3}\), after the definition
\[
\beta_{\mathcal M}
 =
\sup\{t\ge 0:(\mathcal M,d)\Lip(\mathcal M,d^{1-t})\}.
\]
Here \((\mathcal M_1,d_1)\Lip(\mathcal M_2,d_2)\) means that there is a
one-to-one map \(f:\mathcal M_1\to\mathcal M_2\) and constants \(A,B>0\) such
that
\[
A^{-1}d_1(x,y)\le d_2(f(x),f(y))\le B d_1(x,y)
\qquad (x,y\in\mathcal M_1).
\]
Question \(\mathrm{q3}\) asks for a description of the complement of
\(\mathsf{NS_T}\) and, in particular, whether there are non-ultrametric metric
spaces with \(\beta_{\mathcal M}>0\).

\section*{Result}

The result below answers the existential part of Question \(\mathrm{q3}\) on
the literal reading of the question.  It also records a simple family contained
in \(\complement \NS\).

\section*{Proof Intuition}

The obstruction in self-snowflake embeddings is that the snowflake metric
\(d^s\), \(0<s<1\), changes small and large scales at different rates.  A
bounded uniformly discrete metric space has no arbitrarily small or arbitrarily
large nonzero distances.  On the compact interval of possible distances
\([a,b]\), the two functions \(r\) and \(r^s\) are comparable by fixed
constants.  Therefore the identity map is already a Lipschitz embedding from
\((M,d)\) into \((M,d^s)\).  Choosing a bounded uniformly discrete metric that
violates the strong triangle inequality gives the requested non-ultrametric
example.

\begin{lemma}[bounded uniformly discrete spaces self-embed into snowflakes]
\label{lem:bud}
Let \((M,d)\) be a metric space for which there are constants
\(0<a\le b<\infty\) such that
\[
a\le d(x,y)\le b\qquad\text{whenever }x\ne y.
\]
Then for every \(0<s<1\), the identity map is a Lipschitz embedding
\[
(M,d)\Lip(M,d^s).
\]
Consequently, \(\beta_M>0\).  More precisely, for every \(0<t<1\),
\((M,d)\Lip(M,d^{1-t})\), so the supremum is at least \(1\) when
\(\beta_M\) is interpreted on the usual snowflake range \(0<1-t<1\).
\end{lemma}

\begin{proof}
Fix \(0<s<1\) and put \(r=d(x,y)\).  The case \(x=y\) is immediate.  For
\(x\ne y\), \(r\in[a,b]\).  Since \(s-1<0\),
\[
b^{s-1}\le r^{s-1}\le a^{s-1}.
\]
Multiplying by \(r\) gives
\[
b^{s-1}d(x,y)\le d(x,y)^s\le a^{s-1}d(x,y).
\]
Equivalently,
\[
(b^{1-s})^{-1}d(x,y)\le d(x,y)^s\le a^{s-1}d(x,y).
\]
Thus the identity map satisfies the source paper's Lipschitz embedding
inequalities with \(A=b^{1-s}\) and \(B=a^{s-1}\).  Taking \(s=1-t\) proves the
claim about \(\beta_M\).
\end{proof}

\begin{proposition}[explicit non-ultrametric example]
\label{prop:example}
There is an infinite metric space \((M,d)\) that is not an ultrametric and has
\(\beta_M>0\).
\end{proposition}

\begin{proof}
Let
\[
M=\{a,b,c\}\cup\{x_n:n\in\mathbb N\}.
\]
Define \(d(x,x)=0\), make \(d\) symmetric, and set
\[
d(a,b)=1,\qquad d(b,c)=1,\qquad d(a,c)=\frac32,
\]
\[
d(x_n,u)=2\quad(u\in\{a,b,c\},\,n\in\mathbb N),
\qquad
d(x_m,x_n)=1\quad(m\ne n).
\]
We first check that this is a metric.  The only triangle entirely inside
\(\{a,b,c\}\) that is not immediate is
\[
d(a,c)=\frac32\le d(a,b)+d(b,c)=2.
\]
If a triangle has one point \(x_n\) and two points in \(\{a,b,c\}\), the
largest possible side is \(2\), while the sum of the other two sides is at
least \(2+1\); the side joining two core points is at most \(3/2\le 2+2\).
If a triangle has two points \(x_m,x_n\) and one core point, the only
nontrivial inequality is \(2\le 1+2\).  Triangles involving only the \(x_n\)'s
are equilateral.  Hence \(d\) is a metric.

All distinct distances lie in \([1,2]\), so Lemma~\ref{lem:bud} applies and
\(\beta_M>0\).  However, \(M\) is not ultrametric, because
\[
d(a,c)=\frac32>\max\{d(a,b),d(b,c)\}=1.
\]
This gives the requested non-ultrametric example.
\end{proof}

\section*{Verification Notes and Limitations}

The proof is purely metric and uses no external embedding theorem.  It depends
only on the comparability of \(r\) and \(r^s\) on the interval of nonzero
distances of a bounded uniformly discrete metric space.

This packet should be read as a full solution to the yes/no clause of Question
\(\mathrm{q3}\).  It is not a full classification of
\(\complement\mathsf{NS_T}\).  If the intended problem was meant only for
unbounded metric spaces, Banach-space metrics, spaces without isolated points,
or examples not explained by bounded uniform discreteness, this packet does
not settle that stronger interpretation.

\section*{Bounded Novelty Check}

Before promotion, the following bounded checks were performed on 2026-07-18.
The local run indexes were searched:
\begin{itemize}[leftmargin=2em]
\item \texttt{registry\_index.tsv};
\item \texttt{solutions/index.tsv};
\item \texttt{attempts/index.tsv};
\item \texttt{proof\_gaps/index.tsv}.
\end{itemize}
The search terms were:
\begin{itemize}[leftmargin=2em]
\item \texttt{1206.3774}, \texttt{beta\_M}, and \(\beta_{\mathcal M}\);
\item \texttt{other than ultrametrics};
\item \texttt{bounded uniformly discrete} and close variants.
\end{itemize}
The local source and deterministic scan were also searched.  No duplicate
packet or prior run answer to this exact observation was found.

External web searches were also run for phrases around \(\beta_M\), snowflaked
self-embeddings, ultrametrics, and bounded uniformly discrete spaces.  These
found the source arXiv record and nearby snowflake-characterization literature,
but no exact recorded answer to this literal existential clause.  Novelty
confidence is therefore moderate rather than absolute, mainly because the
observation is elementary and may have been considered implicit.

\section*{Human Review Recommendation}

Human review should focus on whether Question \(\mathrm{q3}\) was intended to
exclude bounded uniformly discrete spaces implicitly.  On the literal
definition printed in the source, the proof above gives a valid non-ultrametric
example with \(\beta_M>0\).

\begin{thebibliography}{9}
\bibitem{AlbiacBaudier2012}
Fernando Albiac and Florent Baudier,
\emph{Embeddability of snowflaked metrics with applications to the nonlinear
geometry of the spaces \(L_p\) and \(\ell_p\) for \(0<p<\infty\)},
arXiv:1206.3774, Journal of Geometric Analysis 25 (2015), no. 1, 1--24.
\end{thebibliography}

\end{document}
